\documentclass[11pt]{article}
\usepackage[letterpaper,margin=0.78in]{geometry}
\usepackage{amsmath,amssymb,amsthm,booktabs,microtype}
\usepackage[colorlinks=true,linkcolor=blue!55!black,urlcolor=blue!55!black,citecolor=blue!55!black]{hyperref}
\usepackage{xcolor}
\setlength{\parindent}{0pt}
\setlength{\parskip}{5pt}
\newtheorem*{theorem}{Literature answer}
\newcommand{\Om}{\Omega}
\begin{document}

\begin{center}
{\Large\bfseries Dimension dependence in Hermite--Hadamard inequalities}\\[4pt]
{\large Literature answer to questions in arXiv:1905.03216}\\[5pt]
{\small fa\_banach\_001 $\cdot$ lane 7 $\cdot$ 17 August 2026}
\end{center}

\vspace{3pt}
\noindent\colorbox{blue!7}{\parbox{0.965\linewidth}{
\textbf{Disposition.} Both qualitative dimension questions are already
answered negatively by the direct follow-up arXiv:1907.06122.  No new
mathematical result is claimed.}}

\section*{1. Exact questions}
For convex bodies $\Om\subset\mathbb R^n$, the source first considers the
normalized inequality
\begin{equation}\label{normalized}
 \frac1{|\Om|}\int_\Om f\,dx
 \leq \frac{A_n}{|\partial\Om|}\int_{\partial\Om}f\,d\sigma
\end{equation}
for admissible convex functions.  Immediately after its Theorem~1, it asks
whether the optimal $A_n$ remain bounded with dimension.  Its scale-invariant
theorem concerns positive subharmonic functions and constants $B_n$ in
\begin{equation}\label{scaled}
 \int_\Om f\,dx
 \leq B_n |\Om|^{1/n}\int_{\partial\Om}f\,d\sigma;
\end{equation}
the same paragraph asks whether $B_n$ might decay as $n\to\infty$ \cite{LS}.

The queue's literal phrase ``we do not know'' is not an open problem: it only
says that one step of the Brownian-motion proof does not track the positions of
particles surviving to a fixed time.

\section*{2. Direct literature answer}
\begin{theorem}
The direct follow-up \cite{BBBHLLSS} proves that the optimal constants in
\eqref{normalized} satisfy
\[
 A_n\geq \max\{n-1,1\}.
\]
Consequently no dimension-free constant exists.  It also identifies, on a
fixed convex domain, the optimal constant in
$\int_\Om f\leq c(\Om)\int_{\partial\Om}f$ over positive subharmonic $f$ as
\[
 c(\Om)=\max_{x\in\partial\Om}\partial_{\nu}u(x),
 \qquad -\Delta u=1,\quad u|_{\partial\Om}=0,
\]
where $\nu$ is the inward normal.
\end{theorem}
The lower-bound construction for $A_n$ uses positive convex functions and
therefore lies inside the class asked about in the source.  The fixed-domain
characterization follows by Green's identity for the upper bound and by
Poisson extensions concentrating at a maximizer for the reverse inequality.

\section*{3. Certified non-decay of the scale-invariant constants}
The ellipsoid family displayed in the source already gives the qualitative
answer to the second question, after one normalization correction.  Put
\[
 E_n=\left\{x:\frac{x_1^2}{4}+
 \sum_{j=2}^n\frac{x_j^2}{2n-2}<1\right\},\qquad
 q_n=1-\frac{x_1^2}{4}-\sum_{j=2}^n\frac{x_j^2}{2n-2}.
\]
Direct differentiation gives $-\Delta q_n=3/2$, although the source calls
$q_n$ the torsion function.  Hence the correctly normalized torsion function
is $u_n=(2/3)q_n$.  At $p=(2,0,\ldots,0)$,
\[
 \partial_\nu u_n(p)=\frac23.
\]
Moreover,
\[
 |E_n|^{1/n}
 =\left[\omega_n\,2(2n-2)^{(n-1)/2}\right]^{1/n}
 \longrightarrow 2\sqrt{\pi e}.
\]
The fixed-domain characterization therefore yields
\[
 \liminf_{n\to\infty}B_n\geq \frac{1}{3\sqrt{\pi e}}>0.
\]
Thus the optimal constants in \eqref{scaled} do not decay to zero.  The
follow-up explicitly records the same qualitative conclusion and improves the
upper bound to
$B_n\leq(\omega_n^{1/n}\sqrt n)^{-1}\to(2\pi e)^{-1/2}$.

\section*{4. Verification and scope}
The cited theorem and fixed-domain proposition were checked in the local
source of the follow-up.  The Laplacian, normal derivative, ellipsoid volume,
and Stirling limit above were recomputed independently.  The
normalization slip only changes the numerical lower constant; it does not
alter non-decay.  This packet is a literature audit, not a novel proof packet.

\begin{thebibliography}{9}
\bibitem{LS} J. Lu and S. Steinerberger,
\emph{A Dimension-Free Hermite--Hadamard Inequality via Gradient Estimates for
the Torsion Function}, arXiv:1905.03216 (2019).
\bibitem{BBBHLLSS} T. Beck et al.,
\emph{Improved Bounds for Hermite--Hadamard Inequalities in Higher Dimensions},
arXiv:1907.06122; J. Geom. Anal. 31 (2021), 801--816.
\end{thebibliography}

\end{document}
